\chapter{Statistical Significance Analysis}
\label{app:statistical_significance}

This appendix provides the statistical significance evaluation of the weight estimation results across different sensor modality configurations (ALL, EMG-only, and IMU-only) and validation strategies.

\section{Methodology}
To assess whether the performance differences between sensor configurations are statistically meaningful, we perform the non-parametric Friedman test on the participant-level error metrics ($N = 17$ subjects). The Friedman test is a robust alternative to Repeated Measures ANOVA that does not assume normality of the underlying distributions or sphericity of the differences, making it highly suitable for cross-subject performance comparisons where individual variations frequently violate parametric assumptions.

The Friedman test evaluates the null hypothesis that the ranks of the performances of all modalities across the participants are equal. The test statistic $\chi^2(2)$ is computed based on the sum of ranks for each modality, with degrees of freedom $dof = 2$ for the three compared configurations.

For cases where the Friedman test indicates a significant difference ($p < 0.05$), we perform post-hoc pairwise comparisons using two-tailed Wilcoxon signed-rank tests. To control the family-wise error rate and avoid Type I errors (false positives) arising from multiple comparisons, we apply the Bonferroni correction:
\begin{equation}
    p_{\text{adj}} = \min(1.0, p \times c)
\end{equation}
where $c = 3$ is the number of pairwise comparisons (SF vs. EMG, SF vs. IMU, and EMG vs. IMU).

\section{Friedman and Post-Hoc Results}
Table~\ref{tab:anova_results} summarizes the mean errors, Friedman test statistics, and post-hoc pairwise Wilcoxon comparisons for both the Participant-Specific (5-Fold CV) and Participant-Independent (LOPO CV) settings.

\begin{table*}[h]
    \centering
    \small
    \caption{Friedman test and Bonferroni-corrected post-hoc Wilcoxon signed-rank comparisons of the modalities ($N=17$ participants), computed on per-participant \emph{class-macro} (class-balanced) MAE/RMSE. SF = sensor-fused, EMG = EMG-only, IMU = IMU-only.}
    \label{tab:anova_results}
    \begin{tabularx}{\textwidth}{l c c c c c >{\centering\arraybackslash}X >{\centering\arraybackslash}X >{\centering\arraybackslash}X >{\centering\arraybackslash}p{2.6cm}}
        \toprule
        & \multicolumn{3}{c}{\textbf{Modality Mean MAE/RMSE (kg)}} & \multicolumn{2}{c}{\textbf{Friedman Test}} & \multicolumn{3}{c}{\textbf{Post-hoc $p_{\text{adj}}$}} & \textbf{Pairwise} \\
        \cmidrule(lr){2-4} \cmidrule(lr){5-6} \cmidrule(lr){7-9} \cmidrule(lr){10-10}
        \textbf{Metric \& Strategy} & SF & EMG & IMU & $\chi^2(2)$ & $p$ & SF--EMG & SF--IMU & EMG--IMU & relations \\
        \midrule
        \textbf{MAE (kg)} \\
        \quad Participant-Specific & 0.316 & 0.380 & 0.702 & 34.00 & $<0.001$ & $4.58\times10^{-5}$ & $4.58\times10^{-5}$ & $4.58\times10^{-5}$ & SF $<$ EMG $<$ IMU \\
        \quad Generalized (LOPO)      & 0.536 & 0.546 & 0.989 & 20.59 & $<0.001$ & $1.000$            & $4.58\times10^{-5}$ & $2.29\times10^{-4}$ & SF $\approx$ EMG $<$ IMU \\
        \midrule
        \textbf{RMSE (kg)} \\
        \quad Participant-Specific & 0.441 & 0.520 & 0.974 & 34.00 & $<0.001$ & $4.58\times10^{-5}$ & $4.58\times10^{-5}$ & $4.58\times10^{-5}$ & SF $<$ EMG $<$ IMU \\
        \quad Generalized (LOPO)      & 0.680 & 0.703 & 1.289 & 23.06 & $<0.001$ & $1.000$            & $4.58\times10^{-5}$ & $9.16\times10^{-5}$ & SF $\approx$ EMG $<$ IMU \\
        \bottomrule
    \end{tabularx}
\end{table*}

\section{Discussion of Key Findings}
The Friedman and post-hoc Wilcoxon test results support two primary conclusions:
\begin{enumerate}
    \item \textbf{Sensor Fusion is Strategy-Dependent:} 
    In the Participant-Specific setting, combining EMG and IMU sensors (\texttt{SF}) yields the best performance, outperforming \texttt{EMG-only} ($p_{\text{adj}} = 4.58 \times 10^{-5}$ for MAE) and \texttt{IMU-only} ($p_{\text{adj}} = 4.58 \times 10^{-5}$ for MAE) with high statistical significance.
    Conversely, in the Participant-Independent (LOPO) setting, the multi-modality model does \textit{not} provide a statistically significant benefit over using EMG alone ($p_{\text{adj}} = 1.000$ for both MAE and RMSE). Although the SF model achieves a slightly lower nominal mean MAE (0.536 kg vs. 0.546 kg), this difference is not statistically significant.
    
    \item \textbf{Generalizability Limitations of Kinematics:}
    The \texttt{IMU-only} configuration performs significantly worse than both \texttt{SF} and \texttt{EMG-only} across both evaluation strategies ($p_{\text{adj}} < 0.001$). Crucially, in the LOPO setting, the IMU-only performance degrades severely (MAE = 0.989 kg), indicating that kinematic signatures are highly subject-specific and fail to generalize to unseen users. EMG signals, reflecting direct physiological muscle activation, remain robust and generalize significantly better across individuals.
\end{enumerate}

\section{Model Architecture Significance Analysis}
To evaluate whether the performance differences between different model architectures are statistically meaningful, we perform the non-parametric Friedman test on the participant-level errors ($N=17$) under the Leave-One-Participant-Out (LOPO) CV validation strategy. The architectures compared include the Spatio-Temporal Transformer (ST-Transformer), LSTM, GRU, Naive-Transformer, CNN-GRU, and CNN-LSTM.

The Friedman test results reveal highly significant differences between the architectures for both MAE ($\chi^2(5) = 32.09, p < 0.001$) and RMSE ($\chi^2(5) = 24.56, p < 0.001$). Table~\ref{tab:model_comparison_lopo} presents the mean metrics along with the post-hoc pairwise comparisons using Bonferroni-corrected Wilcoxon signed-rank tests.

\begin{table}[h]
    \centering
    \small
    \caption{Model Architecture Performance and Statistical Significance (LOPO CV, EMG + IMU Modality). Within each column, mean values sharing a common superscript letter are not statistically significantly different at the Bonferroni-corrected level of $\alpha = 0.05$ using Wilcoxon post-hoc comparisons.}
    \label{tab:model_comparison_lopo}
    \begin{tabularx}{\columnwidth}{l >{\centering\arraybackslash}X >{\centering\arraybackslash}X >{\centering\arraybackslash}X}
        \toprule
        Model Architecture & MAE (kg) & RMSE (kg) & $\text{R}^2$ \\
        \midrule
        LSTM & $0.489^{\text{a}}$ & $0.669^{\text{a,b}}$ & $0.846$ \\
        GRU & $0.489^{\text{a}}$ & $0.636^{\text{a}}$ & $0.857$ \\
        ST-Transformer & $0.536^{\text{a,b}}$ & $0.680^{\text{a,b,c}}$ & $0.853$ \\
        Naive-Transformer & $0.622^{\text{b,c}}$ & $0.803^{\text{c}}$ & $0.782$ \\
        CNN-GRU & $0.756^{\text{b,c}}$ & $0.866^{\text{b,c}}$ & $0.706$ \\
        CNN-LSTM & $0.812^{\text{c}}$ & $0.932^{\text{c}}$ & $0.667$ \\
        \bottomrule
    \end{tabularx}
\end{table}
